% =====================================================================
% preamble.tex
% Konfigurasi global untuk Laporan Tugas Akhir IF ITB
% Penulis: Mohammad Nugraha Eka Prawira (NIM 13522001)
% =====================================================================

% --- Encoding dan bahasa ---
\usepackage[T1]{fontenc}
\usepackage[utf8]{inputenc}
% Catatan: babel-indonesian mungkin perlu diinstall pada beberapa distribusi TeX Live.
% Jika tersedia, ganti opsi babel di bawah dengan [english,indonesian] dan
% Catatan: babel-indonesian diaktifkan agar daftar gambar, tabel,
% dan translasi sitasi (seperti dkk.) bekerja secara otomatis.
\usepackage[english,indonesian]{babel}

% --- Judul daftar dengan huruf kapital penuh (pedoman ITB: judul utama BESAR + BOLD) ---
\addto\captionsindonesian{%
  \renewcommand{\contentsname}{DAFTAR ISI}%
  \renewcommand{\listfigurename}{DAFTAR GAMBAR}%
  \renewcommand{\listtablename}{DAFTAR TABEL}%
}

% --- Font ---
\usepackage{mathptmx}        % Times New Roman untuk teks dan matematika
\usepackage[scaled=0.9]{helvet}
\usepackage{courier}

% --- Geometri halaman (Pedoman ITB: top 3cm, bottom 3cm, left 4cm, right 3cm) ---
\usepackage[a4paper,top=3cm,bottom=3cm,left=4cm,right=3cm]{geometry}

% --- Spasi dan paragraf ---
\usepackage{setspace}
% \usepackage{indentfirst} % Dihapus karena tidak ingin ada indent
\setlength{\parindent}{0pt} % Tidak ada indent pada awal paragraf
\setlength{\parskip}{0.8em} % Jarak antar paragraf agar batas paragraf tetap terlihat

% --- Heading customization (ITB style) ---
\usepackage{titlesec}
\renewcommand{\thechapter}{\Roman{chapter}}
\renewcommand{\thesection}{\thechapter.\arabic{section}}
\renewcommand{\thesubsection}{\thesection.\arabic{subsection}}
\renewcommand{\thesubsubsection}{\thesubsection.\arabic{subsubsection}}

% Pedoman ITB: judul utama (judul bab) ukuran 14pt, BESAR + BOLD
\titleformat{\chapter}[display]
  {\normalfont\bfseries\centering}
  {\fontsize{14}{17}\selectfont BAB \thechapter}
  {6pt}
  {\fontsize{14}{17}\selectfont}
\titlespacing*{\chapter}{0pt}{*-2}{20pt}

\titleformat{\section}
  {\normalfont\fontsize{13}{16}\selectfont\bfseries}
  {\thesection}{1em}{}
\titleformat{\subsection}
  {\normalfont\normalsize\bfseries}
  {\thesubsection}{1em}{}
\titleformat{\subsubsection}
  {\normalfont\normalsize\bfseries\itshape}
  {\thesubsubsection}{1em}{}

% --- Daftar isi, gambar, tabel ---
\usepackage[titles]{tocloft}
\usepackage{tocbibind}

\setcounter{tocdepth}{3}
\setcounter{secnumdepth}{3}

% DAFTAR ISI: prefix "BAB" untuk entri chapter (diaktifkan via \addtocontents di main.tex)
\setlength{\cftchapnumwidth}{4em}
% Aktifkan dot leader untuk entri chapter (DAFTAR ISI, BAB I, dst.)
\renewcommand{\cftchapleader}{\cftdotfill{\cftchapdotsep}}
\renewcommand{\cftchapdotsep}{\cftdotsep}

% Perlebar numwidth tiap level agar "III.10" – "III.13.3" tidak overflow
\setlength{\cftsecindent}{1.5em}
\setlength{\cftsecnumwidth}{3.5em}
% Subsection: "III.10.1" perlu ±4.5em; indent = sec_indent + sec_numwidth
\setlength{\cftsubsecindent}{5em}
\setlength{\cftsubsecnumwidth}{4.5em}
% Subsubsection: tetap lebih masuk dari subsection, tetapi tidak terlalu menjorok.
\setlength{\cftsubsubsecindent}{6.8em}
\setlength{\cftsubsubsecnumwidth}{5em}

% DAFTAR GAMBAR: prefix "Gambar" sebelum nomor (misal: "Gambar III.1")
\setlength{\cftfigindent}{0em}
\renewcommand{\cftfigpresnum}{Gambar }
\setlength{\cftfignumwidth}{6.6em}

% DAFTAR TABEL: prefix "Tabel" sebelum nomor (misal: "Tabel III.1")
\setlength{\cfttabindent}{0em}
\renewcommand{\cfttabpresnum}{Tabel }
\setlength{\cfttabnumwidth}{5.6em}

% DAFTAR LAMPIRAN: otomatis via \listoflampiran (baca dari file .loa)
\newlistof{lampiran}{loa}{DAFTAR LAMPIRAN}
\setlength{\cftlampiranindent}{0em}
\setlength{\cftlampirannumwidth}{6.8em}
\renewcommand{\cftlampiranpresnum}{Lampiran }
\renewcommand{\cftlampiranaftersnum}{.}
\renewcommand{\cftlampiranfont}{\bfseries}
\renewcommand{\cftlampiranpagefont}{\bfseries}
\renewcommand{\cftlampiranleader}{\cftdotfill{\cftlampirandotsep}}
\renewcommand{\cftlampirandotsep}{\cftdotsep}

% DAFTAR LAMPIRAN level-2: subbab di dalam tiap Lampiran (\section di lampiran.tex)
\newlistentry[lampiran]{lampiransub}{loa}{1}
\setcounter{loadepth}{2}
\setlength{\cftlampiransubindent}{1.8em}
\setlength{\cftlampiransubnumwidth}{3em}
\renewcommand{\cftlampiransubfont}{\normalfont}
\renewcommand{\cftlampiransubpagefont}{\normalfont}
\renewcommand{\cftlampiransubleader}{\cftdotfill{\cftlampiransubdotsep}}
\renewcommand{\cftlampiransubdotsep}{\cftdotsep}
\renewcommand{\cftlampiransubaftersnumb}{}

% --- Gambar dan tabel ---
\usepackage{graphicx}
\graphicspath{{figures/}}
\usepackage{tikz}
\usetikzlibrary{arrows.meta,positioning,fit,backgrounds,patterns}
\usepackage{pgfplots}
\pgfplotsset{compat=1.18}
\usepackage{caption}
\usepackage{pdflscape}
\captionsetup{labelsep=period,font=small}

% Penomoran eksplisit dengan format Roman (Bab)
\renewcommand{\thefigure}{\thechapter.\arabic{figure}}
\renewcommand{\thetable}{\thechapter.\arabic{table}}
\renewcommand{\theequation}{\thechapter.\arabic{equation}}

\usepackage{booktabs}
\usepackage{longtable}
\usepackage{array}
\usepackage{multirow}
\usepackage{tabularx}
\usepackage{fvextra}

% Kolom tabel bergaris penuh menyerap ruang lebih banyak dari garis vertikal;
% perkecil padding sel agar kolom sempit (mis. kamus atribut node) tidak overfull.
\setlength{\tabcolsep}{4pt}

% --- Daftar terurut ---
\usepackage{enumitem}
\setlist[itemize]{leftmargin=*,topsep=2pt,itemsep=0pt}
\setlist[enumerate]{leftmargin=*,topsep=2pt,itemsep=0pt}

% --- Hyperlink dan referensi ---
\usepackage[
  colorlinks=true,
  linkcolor=black,
  citecolor=black,
  urlcolor=black,
  pdftitle={Pengembangan Arsitektur Perangkat Lunak Companionship Chatbot},
  pdfauthor={Mohammad Nugraha Eka Prawira}
]{hyperref}

% --- Bibliografi dengan natbib + plainnat ---
\usepackage[round,authoryear,sort&compress]{natbib}
\bibliographystyle{bibliography/itbnat}
\let\textcite\citet
\let\parencite\citep

% --- Utilities ---
\usepackage[table]{xcolor}
\usepackage{microtype}

% Warna abu-abu untuk baris header tabel, mengikuti contoh pada pedoman penulisan TA
\definecolor{tableheadgray}{gray}{0.8}
% Toleransi akhir untuk paragraf yang memuat literal \texttt panjang (mengurangi overfull)
\setlength{\emergencystretch}{3em}
% Izinkan pemenggalan baris setelah underscore pada identifier kode (mis. \texttt{nama_panjang})
\makeatletter
\let\@origunderscore\_
\renewcommand{\_}{\@origunderscore\allowbreak}
\makeatother
\usepackage{csquotes}
\DeclareQuoteStyle{indonesian}
  {\quotedblleft}{\quotedblright}
  {\quoteleft}{\quoteright}

% --- Header dan footer kosong (ITB standar) ---
\usepackage{fancyhdr}
\pagestyle{plain}

% --- Header chapter tidak boleh kosong di halaman pertama bab ---
\makeatletter
\let\ps@plain\ps@plain
\makeatother
